\scanpage{72}
\uline{Remarque} Dans deux contextes différents
``déploiement''$^{*}$ de \deleted{\unsure{strag}} stratifications, on
a rencontré des espaces topologiques ordonnés\inserted{$\Sigma$},
satisfaisant les conditions suivantes

\margin{$*$ en deux sens assez différents}

a) $\Sigma_{(1)}\xrightarrow{\uline{s}}\Sigma_{\deleted{\unsure{*|0}}}=\Sigma_{(0)}$
est une fibration propre

b) $\Sigma_{(1)}\xrightarrow{\uline{b}}\Sigma_{(0)}$ est une immersion propre

c) Si $x,y\in\Sigma$ \unsure{\deleted{sont}|sont} \unsure{ds|dans} une $\unsure{\overline{\mathrm m}}$ composante connexe,
alors $x\leq y\Longrightarrow x=y$
\inserted{(i.e. si $x,y\in\Sigma$, $x<y$, alors $x,y$ sont dans deux
composantes connexes différentes)}

\margin{conséquence de a) b) si on suppose que
$\Sigma(1)^{*}\xrightarrow{\uline b}\Sigma$ est \uline{\unsure{rare}},
et que les comp. connexes de $\Sigma$ de dim finie\ldots}

Dans un cas (celui d'une notion assez \unsure{anodine|analogue à}
\unsure{un contexte|un certain|de la notion|la notion} de déploiement d'une stratification de type $I$ disons,
par $\Sigma=\coprod X_i$,
$\Sigma_1=\displaystyle\coprod_{\substack{i,j\text{ tels que}\\X_i\subset X_j}}X_i$),
$\uline s$ est étale, (\unsure{localt} fini donc) et rien de plus spécial pour $\uline b$.
Dans le cas de la théorie \deleted{\unsure{proposée}}\inserted{développée}
ici des déploiements\inserted{de stratifications équisingulières},
$\uline b$ est bordant, rien de spécial sur $\uline s$.
Dans l'un et l'autre cas, on peut appliquer \unsure{le|la} \unsure{sorite|sortie}
\unsure{du|des} $\widetilde I=\pi_0(\Sigma)$ --- dans le premier cas,
cela revient à remplacer la stratification par les $X_i$ par celle
par les comp. connexes des $X_i$ (donnant un $\widetilde I$;
dans le deuxième,
\deleted{à \unsure{remplacer|raffiner} \unsure{ito} $\unsure{\Longleftarrow}$ \unsure{exact}}
itou pour les
\[
X_i^{*}=X_i\setminus\bigcup_{j<i}\unsure{X_{ij}|X_j}.
\]
\mspage{71}
